\chapter{Gynaecomastia}
\index{Gynaecomastia}

\section*{Definition and Overview}
Gynaecomastia is the benign enlargement of breast tissue in males, caused by an imbalance between the effects of oestrogen and testosterone. It is a common condition that can affect newborns, boys during puberty, and older men, and it is usually harmless and temporary. Gynaecomastia appears as a firm, tender disc of tissue beneath the nipple, often on both sides, and must be distinguished from fat deposition (pseudogynaecomastia) and, rarely, breast cancer. Most cases resolve on their own or with treatment of an underlying cause, and surgery is an option for persistent or troublesome enlargement. Reassurance is often the most important part of management.

\section*{Epidemiology}
Gynaecomastia is very common. Up to 60--70\% of boys develop some breast tissue during puberty, usually between ages 13 and 15, and it resolves within one to two years in the majority. It also affects up to 50--65\% of adult men at some stage, with prevalence increasing with age. Around a third of cases have no identifiable cause. Gynaecomastia is more common in men who are overweight, who drink alcohol, who use certain medicines or anabolic steroids, and in those with conditions affecting testosterone levels, such as cirrhosis and kidney disease.

\section*{Aetiology and Pathophysiology}
Gynaecomastia results from an imbalance in the hormones that control male breast tissue.

\begin{itemize}
    \item \textbf{Oestrogen effect:} breast tissue grows when the relative effect of oestrogen exceeds that of testosterone
    \item \textbf{Physiological gynaecomastia:} newborn (maternal oestrogen), puberty (temporary imbalance), and old age (declining testosterone)
    \item \textbf{Medicines:} spironolactone, some anti-ulcer and antipsychotic medicines, digoxin, androgens and anabolic steroids, and some prostate cancer treatments
    \item \textbf{Substance use:} alcohol, marijuana, heroin, and amphetamines
    \item \textbf{Hormonal conditions:} low testosterone, hyperthyroidism, and rare oestrogen-producing tumours
    \item \textbf{Systemic disease:} liver cirrhosis, kidney failure, and malnutrition alter hormone balance
    \item \textbf{Genetic conditions:} Klinefelter syndrome and other disorders
    \item \textbf{Idiopathic:} in many cases, no cause is found
\end{itemize}

\section*{Clinical Features}
Gynaecomastia has characteristic findings on examination.

\begin{itemize}
    \item A firm, rubbery, disc-like mass beneath the nipple
    \item Bilateral involvement in most cases, though it can be one-sided
    \item Tenderness or pain, which is common during puberty and early growth
    \item Breast enlargement that may be gradual or rapid
    \item No discharge from the nipple in simple gynaecomastia
    \item Symptoms of an underlying cause when present, such as reduced libido, or weight gain
    \item A normal appearance of the testicles and genitalia
\end{itemize}

\section*{Differential Diagnosis}
Other causes of male breast enlargement and breast lumps must be considered.

\begin{itemize}
    \item Pseudogynaecomastia: enlargement due to fat rather than glandular tissue, common in overweight men
    \item Male breast cancer, which is rare but serious; it tends to be painless, hard, fixed, one-sided, and may cause nipple discharge or skin changes
    \item Lipoma and other benign breast lumps
    \item Cysts and infection of the breast
    \item Nipple and skin lesions
    \item Testicular tumours, which can increase oestrogen and cause gynaecomastia
    \item Medications and substance-related causes
\end{itemize}

\section*{When to Seek Medical Care}
See a doctor for any new breast lump or enlargement in a male, particularly if it is one-sided, hard, fixed, growing rapidly, or associated with nipple discharge or skin changes. Persistent or painful gynaecomastia also warrants review.

\textbf{Contact your local emergency services for:}
\begin{itemize}
    \item This condition is not an emergency; however, any rapidly growing breast mass with chest pain or difficulty breathing needs urgent assessment
    \item Severe pain, redness, or signs of infection in the breast
\end{itemize}

\section*{Diagnosis and Investigations}
The cause of gynaecomastia is determined by history, examination, and targeted tests.

\begin{itemize}
    \item \textbf{History:} age, medicines, alcohol and drug use, and symptoms of hormone imbalance
    \item \textbf{Examination:} confirming glandular tissue beneath the nipple, and checking the testicles
    \item \textbf{Blood tests:} testosterone, oestrogen, thyroid function, liver function, and kidney function
    \item \textbf{Breast ultrasound or mammogram:} if cancer is suspected or the diagnosis is uncertain
    \item \textbf{Imaging of the testicles:} ultrasound if a testicular cause is suspected
    \item \textbf{Review of medicines:} identifying drugs that cause gynaecomastia
    \item \textbf{Reassurance and monitoring:} in most cases, no further tests are needed
\end{itemize}

\section*{Management}
Management depends on the cause, severity, and how much the gynaecomastia bothers the person.

\begin{itemize}
    \item \textbf{Reassurance and observation:} most pubertal gynaecomastia resolves on its own within two years
    \item \textbf{Treating the cause:} stopping or changing culprit medicines, reducing alcohol, and treating underlying conditions
    \item \textbf{Weight management:} reduces fat and the appearance of enlargement, though glandular tissue persists
    \item \textbf{Medicines:} tamoxifen or other agents may be used for painful or persistent cases
    \item \textbf{Surgery:} liposuction and glandular excision for persistent, large, or distressing gynaecomastia
    \item \textbf{Follow-up:} monitoring for resolution or progression
    \item \textbf{Psychological support:} addressing the effects on confidence and body image
\end{itemize}

\section*{Complications}
\begin{itemize}
    \item Psychological distress, embarrassment, and effects on body image and quality of life
    \item Persistence of breast tissue beyond puberty
    \item Pain and tenderness
    \item Complications of surgery, including scarring, asymmetry, and altered nipple sensation
    \item Missed diagnosis of the rare male breast cancer if lumps are assumed benign
    \item Effects of underlying hormonal or systemic disease if not identified
\end{itemize}

\section*{Prognosis and Outcomes}
The prognosis of gynaecomastia is excellent. Pubertal gynaecomastia resolves spontaneously in about 80--90\% of boys within one to two years, and gynaecomastia caused by medicines or substances improves when they are stopped. In persistent cases, surgery provides good cosmetic results. Gynaecomastia itself is benign and does not increase the risk of breast cancer, though the underlying breast tissue should be examined if it changes. Most men find that the condition causes little or no long-term problem.

\section*{Prevention and Self-Care}
\begin{itemize}
    \item Avoid anabolic steroids and other substances that cause gynaecomastia
    \item Limit alcohol consumption
    \item Maintain a healthy weight
    \item Review medicines with a doctor if breast tissue develops after starting a new medicine
    \item Perform regular self-examination of the breast area and see a doctor for any new lump
    \item Seek review of one-sided, hard, or growing lumps promptly
    \item Reassure adolescent boys that pubertal gynaecomastia is normal and usually temporary
    \item Seek support if body image concerns are affecting confidence
\end{itemize}

\section*{Sources and Further Reading}
The following reputable sources provide further information for healthcare students and patients seeking reliable, current advice about gynaecomastia.

\begin{itemize}
    \item Healthdirect Australia. \textit{Gynaecomastia.} https://www.healthdirect.gov.au/
    \item Australian Society of Plastic Surgeons. \textit{Gynaecomastia surgery information.} https://plasticsurgery.org.au/
    \item Andrology Australia. \textit{Male breast enlargement information.} https://www.andrologyaustralia.org/
    \item NHS. \textit{Gynaecomastia (male breast enlargement).} https://www.nhs.uk/
    \item Mayo Clinic. \textit{Gynecomastia: symptoms and causes.} https://www.mayoclinic.org/
    \item MSD Manual. \textit{Gynaecomastia.} https://www.msdmanuals.com/
\end{itemize}
